% !TeX program = xelatex
% !TeX encoding = UTF-8


\clearpage
\chapter*{一些数学备注}
\addcontentsline{toc}{chapter}{一些数学备注}

在自动翻译过程中，校者要求ChatGPT对一些数学上的潜在问题进行检查。本节记录这个检查的结果：原书中若干处公式、估计或构造上的疑点，并给出相应的计算或理由。其中一些可能只是笔误；另一些可能需要补充条件，或澄清原文的用语。局部论证存在问题，并不意味着相关定理的结论不成立。下文提出的修正方向也不视为已经完成的替代证明。希望读者在阅读时自行对照是否需要进行勘误。


\section*{推论 3.6 中上下界的选择}
\Source{第 24 页，推论 3.6 的证明}

原文将正侧的 $2\varepsilon''$ 取为最小上界，将负侧的 $2\varepsilon'$ 取为最大下界，随后断言相应水平带包含在 $V$ 内。然而，要排除 $V$ 外的正侧点，需要控制的是正侧的正下界。负侧则相应需要负的上界。

例如，若正侧在 $V$ 外的函数值介于 $0.1$ 与 $10$ 之间，取上确界的一半为 $5$，仍会把部分外部点纳入水平带。因此，仅凭原文所取的上下界，不能推出所需的包含关系。

这里疑似交换了上界与下界。另一种可能是构造函数时还需满足某个未写出的限制。

\section*{定理 3.14 中用于定义收缩的方程无解}
\Source{第 33--34 页，定理 3.14 证明的第一步}

原文用下列方程的正实数解定义 $\rho$：
\[
-\frac{|x|^2}{\rho^2}+\rho^2|y|^2
=(-|x|^2+|y|^2)(1-t)-t.
\]
取 $x=0$、$|y|=1/2$、$t=1$。所取点属于原文定义的 $L_\lambda$，也满足 $|y|\geq1/10$。此时方程化为
\[
\rho^2/4=-1,
\]
因而不存在正实数解。原式不能在所述的全部参数范围内定义收缩。

后续对 $\rho$ 取最大值的操作，不能直接消除这个问题，因为参与取最大值的正解本身可能不存在。这里可能需要单独处理 $x=0$ 的情形，或给出在全域上有定义且连续的截断公式。


\section*{第一消去定理局部模型中的额外平方项}
\Source{第 55 页，第一消去定理证明中的函数 $F$}

原文给出的局部模型为
\[
F(x)=f(p)+2\int_0^{x_1}v(t)\,dt
-x_1^2-\cdots-x_{\lambda+1}^2
+x_{\lambda+2}^2+\cdots+x_n^2.
\]
预备假设要求 $v(t)=t$ 在 $0$ 附近成立。因此，在原点附近有
\[
2\int_0^{x_1}v(t)\,dt-x_1^2=0.
\]
$F$ 在该邻域内不依赖 $x_1$，故其 Hessian 在 $x_1$ 方向退化。这与原文随后的非退化临界点陈述不符。

另设 $e=(1,0,\ldots,0)$。由预备假设 $v(1)=0$，可得
\[
\frac{\partial F}{\partial x_1}(e)=2v(1)-2=-2.
\]
所以，按照印出的公式，$e$ 甚至不是临界点。

一个很可能的笔误是负平方和多写了 $x_1^2$。若将该和从 $x_2^2$ 开始，上述两个直接矛盾便可消除。原文随后给出的 $F(e)$ 的积分表达式也支持这一判断。不过，完整修正仍需与后续构造一起核对。

\section*{引理 5.9 的严格不等式遗漏原点}
\Source{第 65 页，引理 5.9 的陈述}

原文对相应邻域中的点断言
\[
|\partial_t h_t(x)|<K|x|,
\qquad
|\pi_a h_t(x)|>k|x|\quad(x\in\R^a).
\]
这里 $h_t(0)=0$ 对每个 $t$ 都成立，故 $\partial_t h_t(0)=0$。将 $x=0$ 代入两个严格不等式，分别得到
\[
0<0,\qquad 0>0.
\]
因此，所述的全称断言在原点不成立。

这应是陈述范围的遗漏：可以明确要求 $x\ne0$，或说明原点处仅成立非严格不等式。这一问题不否定相应估计对非零点成立。


\section*{定理 5.6 证明末尾的半径估计}
\Source{第 65--66 页，定理 5.6 的证明末尾}

设 $E_1$ 是外球 $E$ 中的内球。原文以 $r$ 表示 $E$ 的半径，却对
$x\in\R^a\cap(E\setminus E_1)$ 使用估计
\[
|\pi_a h_{t_0}(x)|>kr.
\]
引理 5.9 直接给出的只是
\[
|\pi_a h_{t_0}(x)|>k|x|.
\]
由于这些点可能满足 $|x|<r$，后一不等式不能推出前一不等式。这里疑似混淆了外球半径与环形区域中 $|x|$ 的下界。

若改用内球半径，可以得到另一种下界。然而，原文还要据此推出归纳过程有限次终止。因此，除了修正这一估计，还需说明所得步长具有统一的正下界。

此外，同页在选择延拓速度时出现未定义的常数 $k_1$，随后又使用 $K$。这些符号之间的关系并不清楚。半径估计、常数的依赖关系以及终止性论证，宜作为同一个问题考察。

\section*{引理 6.13 中平行移动与法向性的关系}
\Source{第 81 页，引理 6.13 的证明}

原文沿曲线 $C$ 平行移动初始标架，并以平行移动保持内积为由，断言所得到的向量始终垂直于 $U'$。但保持两个平行向量之间的内积，并不自动意味着保持与曲线切向量的正交关系。

设 $C=C(t)$，$\tau$ 是 $C$ 的单位切向量，$\xi$ 沿 $C$ 平行。由度量相容性和 $\nabla_{\dot C}\xi=0$，有
\[
\frac{d}{dt}\langle\xi,\tau\rangle
=\langle\xi,\nabla_{\dot C}\tau\rangle.
\]
原文未说明 $C$ 对所构造的度量是测地线，故右端不必为零。初始时垂直于 $C$ 的向量，经过平行移动后未必仍垂直于 $C$。

即使 $M$ 是全测地的，也只能据此保证初始切于 $M$ 的平行向量仍切于 $M$，不能直接推出它仍垂直于 $C$。因此，原文关于法向性的理由似乎缺少一个条件或构造步骤。

可能的处理方式是采用曲线法丛上的投影联络，或进一步选择使 $C$ 成为测地线的度量。无论采用哪种方式，都还需检验它与两端预定标架的相容性。


\section*{两处“横截相交”与维数条件的冲突}
\Source{第 84 页与第 94 页}

按照定义 5.1，子流形 $A,B\subset X$ 在交点 $z$ 横截，意味着
\[
T_zA+T_zB=T_zX.
\]
由此必有 $\dim A+\dim B\geq\dim X$。

第 84 页称 $U'$ 与 $M$ 沿 $C$ 横截相交。此时
\[
\dim U'=2,\qquad \dim M=r,\qquad
\dim V=r+s,\qquad s\geq3.
\]
因此 $2+r<r+s$，二者在有交点时不可能满足上述横截条件。

第 94 页称一维弧分别与 $S_L^{\lambda-1}$ 和 $S_R^{n-\lambda-1}$ 在 $(n-1)$ 维流形 $V_0$ 中横截相交。由 $2\leq\lambda\leq n-2$，两次维数和分别满足
\[
1+(\lambda-1)=\lambda<n-1,
\qquad
1+(n-\lambda-1)=n-\lambda<n-1.
\]
这两处相交同样不可能是定义 5.1 意义下的横截相交。

原文可能在宽泛地使用“横截”一词。例如，第二处实际需要的可能只是弧在交点处不切于球面。第一处则需要另行明确沿 $C$ 的切空间交条件。这属于术语与所需几何条件之间的疑点，不宜仅替换一个词而不检查后续用途。

\section*{指标 1 消去证明中修改向量场的位置}
\Source{第 105 页，指标 1 临界点的消去证明}

原文要求在 $V_{2+}$ 右侧改变类梯度向量场 $\xi$，以便让 $V_{1+}$ 中的球面 $S$ 避开指标 2 临界点的左球面。

然而，该左球面在 $V_{1+}$ 中的位置，由 $V_{1+}$ 与相应指标 2 临界点之间的轨线决定。若只在 $V_{2+}$ 以上改变向量场，而这段轨线所在区域保持不变，则左球面在 $V_{1+}$ 中的位置也不会改变。这样便不能实现原文所要求的避交。

这里疑似把 $V_{1+}$ 误写为 $V_{2+}$。一个可能的修改区域位于 $V_{1+}$ 右侧、相应指标 2 临界点之前。是否确应如此，应结合引理 4.7 中通过同痕改变类梯度向量场的具体构造判断。

 
